\documentclass[11pt]{article}

\usepackage[T1]{fontenc}
\usepackage{lmodern}
\usepackage{microtype}
\usepackage[a4paper,margin=30mm]{geometry}
\usepackage{amsmath,amssymb,amsthm,mathtools}
\usepackage[hidelinks]{hyperref}

\allowdisplaybreaks[4]
\emergencystretch=2em

\newtheorem{theorem}{Theorem}[section]
\newtheorem{lemma}[theorem]{Lemma}
\newtheorem{proposition}[theorem]{Proposition}
\newtheorem{corollary}[theorem]{Corollary}
\theoremstyle{definition}
\newtheorem{definition}[theorem]{Definition}
\theoremstyle{remark}
\newtheorem{remark}[theorem]{Remark}

\newcommand{\N}{\mathbb{N}}
\newcommand{\Nzero}{\mathbb{N}_0}
\newcommand{\Zfortynine}{\mathbb{Z}/49\mathbb{Z}}
\newcommand{\cpl}[1]{#1^{\mathsf{c}}}
\newcommand{\ints}[2]{[#1,#2]_{\mathbb{Z}}}

\title{A Fixed Aperiodic Greedy Pair-Sum-Avoiding Sequence}
\author{Zhiheng Li}
\date{}

\begin{document}

\maketitle

\begin{abstract}
With the help of GPT-5.6 Sol, we construct an explicit finite set of positive integers whose greedy
pair-sum-avoiding extension has a non-eventually periodic sequence of
consecutive gaps.  The construction consists of a nonperiodic scale-eight
sumset controller embedded in three residue classes modulo \(49\), together
with a finite modular shield.  For the fixed \(21\)-element seed
\[
\begin{split}
A=\{&1,2,3,5,7,13,22,27,28,32,36,40,47,48,52,\\
    &63,71,77,81,89,97\},
\end{split}
\]
we give an explicit infinite set \(S\), prove
\(t\in S\Longleftrightarrow t\notin S+S\) for every \(t>97\), and prove
that the membership indicator of \(S\) is not ultimately periodic.  The
exact recurrence verifies the least-admissible-next-term rule, with equal
summands allowed, and implies that the gap sequence of the greedy extension
is not eventually periodic.
\end{abstract}

\section{The greedy extension and the main theorem}

Write \(\N=\{1,2,3,\ldots\}\) and
\(\Nzero=\{0,1,2,\ldots\}\).  For subsets \(X,Y\subseteq\Nzero\) and
\(c\in\Nzero\), put
\[
X+Y=\{x+y:x\in X,\ y\in Y\},
\qquad
c+X=\{c+x:x\in X\},
\qquad
\cpl X=\Nzero\setminus X.
\]
Equal summands are allowed in every sumset.

Let \(A=\{a_1<\cdots<a_k\}\subseteq\N\) be a nonempty finite set.  Install
all elements of \(A\) as the initial terms.  Having chosen
\(a_1<\cdots<a_n\), with \(n\geq k\), define
\begin{equation}\label{eq:greedy-definition}
a_{n+1}
=\min\{t>a_n:t\notin A_n+A_n\},
\qquad
A_n=\{a_1,\ldots,a_n\}.
\end{equation}
The minimum exists because \(A_n+A_n\) is finite.  No admissibility condition
is imposed on the initial set \(A\).  Let
\(d_m=a_{m+1}-a_m\) denote the consecutive gaps.

\begin{theorem}\label{thm:main}
Let
\[
\begin{split}
A=\{&1,2,3,5,7,13,22,27,28,32,36,40,47,48,52,\\
    &63,71,77,81,89,97\}.
\end{split}
\]
For the sequence generated from \(A\) by~\eqref{eq:greedy-definition},
\[
\forall M\geq1\ \forall p\geq1\ \exists m\geq M
\quad d_{m+p}\neq d_m.
\]
Thus its consecutive gap sequence is not eventually periodic.
\end{theorem}

The following elementary lemma will identify an explicitly constructed set
with the output of the least-next rule.

\begin{lemma}[Exact recurrence implies the greedy rule]\label{lem:greedy}
Let \(S\subseteq\N\) be infinite, let \(N\in S\), and put
\(A=S\cap\ints1N\).  Suppose that
\begin{equation}\label{eq:exact-recurrence}
t\in S
\quad\Longleftrightarrow\quad
t\notin S+S
\qquad (t>N).
\end{equation}
Then the increasing enumeration of \(S\) is exactly the greedy extension of
\(A\).
\end{lemma}

\begin{proof}
Regard the least-next rule as scanning the integers \(N+1,N+2,\ldots\) in
increasing order and selecting precisely the admissible ones.  Before the
first decision, the selected set is \(A=S\cap\ints1N\).  We prove by
induction on \(t>N\) that, after the decision at \(t\), the selected set up to
\(t\) is \(S\cap\ints1t\).  Before the decision at \(t\), the induction
hypothesis gives the selected set \(S\cap\ints1{t-1}\).  If \(t=u+v\) with
\(u,v\in S\), then \(u,v\geq1\), and hence \(u,v<t\).  Consequently
\[
t\in(S\cap\ints1{t-1})+(S\cap\ints1{t-1})
\quad\Longleftrightarrow\quad
t\in S+S.
\]
Thus \(t\) is admissible exactly when \(t\notin S+S\), which
by~\eqref{eq:exact-recurrence} is exactly when \(t\in S\).  This proves the
induction.  Equivalently, if \(s_m<s_{m+1}\) are consecutive elements of
\(S\) with \(s_m\geq N\), then
\[
s_{m+1}\notin S_m+S_m,
\qquad
s_m<t<s_{m+1}\Longrightarrow t\in S_m+S_m,
\]
where \(S_m=\{s_1,\ldots,s_m\}\).  Hence \(s_{m+1}\) is the least admissible
integer larger than \(s_m\), as required.
\end{proof}

\section{A scale-eight sumset controller}

For integers \(a\leq b\), write
\(\ints ab=\{a,a+1,\ldots,b\}\).  We shall repeatedly use the elementary
identity
\begin{equation}\label{eq:interval-sum}
\ints ab+\ints cd=\ints{a+c}{b+d}
\qquad (a\leq b,\ c\leq d).
\end{equation}
Indeed, only the reverse inclusion requires comment.  If
\(a+c\leq z\leq b+d\), set \(x=\max\{a,z-d\}\) and \(y=z-x\).  Then
\(a\leq x\leq b\), \(c\leq y\leq d\), and \(z=x+y\).

Define
\begin{equation}\label{eq:L}
L=\bigcup_{k\geq0}\ints{8^k}{4\cdot8^k-1}
\subseteq\Nzero
\end{equation}
and
\begin{align}
Y_0&=\{0\}\cup
      \bigcup_{k\geq0}\ints{4\cdot8^k}{8\cdot8^k-1},
      \label{eq:Y0}\\
Y_1&=\ints13\cup
      \bigcup_{k\geq0}\ints{16\cdot8^k-1}{32\cdot8^k-1},
      \label{eq:Y1}\\
Y_2&=\ints01\cup\ints7{15}\cup
      \bigcup_{k\geq0}\ints{64\cdot8^k-1}{128\cdot8^k-1}.
      \label{eq:Y2}
\end{align}
For every \(k\geq0\), the interval
\(\ints{8^k}{8^{k+1}-1}\) is the disjoint union
\[
\ints{8^k}{4\cdot8^k-1}
\mathbin{\dot\cup}
\ints{4\cdot8^k}{8\cdot8^k-1}.
\]
These intervals partition \(\N\).  Together with \(0\in Y_0\), this gives
\begin{equation}\label{eq:Y0-complement-L}
Y_0=\cpl L.
\end{equation}

\begin{lemma}[Controller identities]\label{lem:controller}
The sets in~\eqref{eq:Y0}--\eqref{eq:Y2} satisfy
\begin{align}
Y_1&=\cpl{(Y_0+Y_0)},                                      \label{eq:gate01}\\
Y_2&=\cpl{(Y_1+Y_1)},                                      \label{eq:gate12}\\
Y_0&=\{0\}\cup\bigl(1+\cpl{(Y_2+Y_2)}\bigr).             \label{eq:gate20}
\end{align}
\end{lemma}

\begin{proof}
We first compute the three self-sumsets.  Put
\(I_k=\ints{4\cdot8^k}{8\cdot8^k-1}\).  For \(a=8^k\),
\[
0+I_k=\ints{4a}{8a-1},
\qquad
I_k+I_k=\ints{8a}{16a-2},
\]
so these adjacent intervals fill \(\ints{4a}{16a-2}\).  If
\(i\leq k\), then \(8^i\leq a\), and~\eqref{eq:interval-sum} gives
\[
I_i+I_k
=\ints{4(8^i+a)}{8(8^i+a)-2}
\subseteq\ints{4a}{16a-2}.
\]
Every sum of two nonzero elements of \(Y_0\) is of this form after
interchanging the summands.  Therefore
\begin{equation}\label{eq:Y0-square}
Y_0+Y_0
=\{0\}\cup\bigcup_{k\geq0}
  \ints{4\cdot8^k}{16\cdot8^k-2}.
\end{equation}

Next put \(C=\ints13\) and
\(J_k=\ints{16\cdot8^k-1}{32\cdot8^k-1}\).  For \(a=8^k\),
\begin{align*}
C+C&=\ints26,\\
C+J_k&=\ints{16a}{32a+2},\\
J_k+J_k&=\ints{32a-2}{64a-2}.
\end{align*}
The latter two intervals overlap and fill \(\ints{16a}{64a-2}\).
If \(i\leq k\), then
\[
J_i+J_k
=\ints{16(8^i+a)-2}{32(8^i+a)-2}
\subseteq\ints{16a}{64a-2},
\]
because \(8^i\leq a\) and \(8^i\geq1\).  Hence
\begin{equation}\label{eq:Y1-square}
Y_1+Y_1
=\ints26\cup\bigcup_{k\geq0}
  \ints{16\cdot8^k}{64\cdot8^k-2}.
\end{equation}

Finally let \(H=\ints01\cup\ints7{15}\) and
\(K_k=\ints{64\cdot8^k-1}{128\cdot8^k-1}\).  Directly
from~\eqref{eq:interval-sum},
\begin{align*}
H+H&=\ints02\cup\ints7{30},\\
(\ints01)+K_k&=\ints{64a-1}{128a},\\
(\ints7{15})+K_k&=\ints{64a+6}{128a+14},\\
K_k+K_k&=\ints{128a-2}{256a-2},
\end{align*}
where again \(a=8^k\).  Since \(64a+6\leq128a\), the two middle
intervals overlap, and hence
\[
H+K_k=\ints{64a-1}{128a+14}.
\]
This interval overlaps \(K_k+K_k\), and together they fill
\(\ints{64a-1}{256a-2}\).  For \(i\leq k\), one also has
\[
K_i+K_k
=\ints{64(8^i+a)-2}{128(8^i+a)-2}
\subseteq\ints{64a-1}{256a-2}.
\]
It follows that
\begin{equation}\label{eq:Y2-square}
Y_2+Y_2
=\ints02\cup\ints7{30}\cup\bigcup_{k\geq0}
  \ints{64\cdot8^k-1}{256\cdot8^k-2}.
\end{equation}

The complementary gaps between the long intervals
in~\eqref{eq:Y0-square} are
\[
\ints13,
\qquad
\ints{16\cdot8^k-1}{32\cdot8^k-1}\qquad(k\geq0),
\]
which proves~\eqref{eq:gate01}.  The complementary gaps
in~\eqref{eq:Y1-square} are
\[
\ints01,
\qquad
\ints7{15},
\qquad
\ints{64\cdot8^k-1}{128\cdot8^k-1}\qquad(k\geq0),
\]
which proves~\eqref{eq:gate12}.

To prove the feedback identity, subtract \(1\) from the blocks of \(L\).
The blocks of \(L-1=\{\ell-1:\ell\in L\}\) corresponding to \(k=0\) and
\(k=1\) are \(\ints02\) and \(\ints7{30}\), respectively.  For \(k=j+2\),
the corresponding block is
\[
\ints{8^{j+2}-1}{4\cdot8^{j+2}-2}
=\ints{64\cdot8^j-1}{256\cdot8^j-2}.
\]
Thus~\eqref{eq:Y2-square} is exactly
\begin{equation}\label{eq:Y2-square-L}
Y_2+Y_2=L-1.
\end{equation}
For \(q\geq1\), equations~\eqref{eq:Y0-complement-L}
and~\eqref{eq:Y2-square-L} give
\[
q\in Y_0
\iff q\notin L
\iff q-1\notin L-1
\iff q-1\in\cpl{(Y_2+Y_2)}.
\]
Together with \(0\in Y_0\), this is~\eqref{eq:gate20}.
\end{proof}

\begin{lemma}\label{lem:Y0-aperiodic}
The membership indicator of \(Y_0\) is not ultimately periodic.
\end{lemma}

\begin{proof}
Suppose that there are \(N\geq0\) and \(h\geq1\) such that
\[
n\in Y_0\iff n+h\in Y_0
\qquad(n\geq N).
\]
Choose \(a=8^k>\max\{N,h\}\) and set \(m=4a-h\).  Then
\[
a\leq m\leq4a-1,
\qquad
m\geq N.
\]
Hence \(m\in L\), so \(m\notin Y_0\), whereas
\(m+h=4a\in\ints{4a}{8a-1}\subseteq Y_0\).  This contradicts the assumed
periodicity.
\end{proof}

\section{A finite shield modulo 49}

All sums in this section are taken in \(\Zfortynine\), with residues written
as their representatives in \(\{0,1,\ldots,48\}\).  Put
\begin{equation}\label{eq:PBD}
P=\{7,14,28\},
\qquad
B=\{3,22,32,40,48\},
\qquad
D=\{1,2,5,13,27,36,47\},
\end{equation}
and let \(Q=B\cup P\).

\begin{lemma}[Finite shield certificate]\label{lem:shield}
In \(\Zfortynine\),
\begin{align}
(B+B)\cup(P+B)\cup(D+B)&=\Zfortynine\setminus Q,
                                                        \label{eq:cover-mod}\\
Q\cap\bigl((B+B)\cup(P+B)\cup(D+B)\cup(D+P)\bigr)
&=\varnothing,                                         \label{eq:avoid-mod}\\
(P+P)\cap Q&=P.                                        \label{eq:PP-mod}
\end{align}
Moreover, the only ordered pairs from \(P\times P\) whose sum lies in \(Q\)
are the three diagonal pairs
\begin{equation}\label{eq:diagonal-pairs}
7+7=14,
\qquad
14+14=28,
\qquad
28+28=49+7.
\end{equation}
\end{lemma}

\begin{proof}
The required finite calculation is displayed in full:
\begin{align*}
B+B={}&\{2,5,6,13,15,21,23,25,31,35,39,43,44,47\},\\
P+B={}&\{1,5,6,10,11,13,17,19,27,29,31,36,39,46,47\},\\
D+B={}&\{0,1,4,5,8,9,10,12,16,18,19,20,23,24,26,27,\\
      &\hspace{2.5em}30,33,34,35,37,38,39,41,42,45,46\},\\
D+P={}&\{1,5,6,8,9,12,15,16,19,20,26,27,29,30,33,34,41,43\},\\
P+P={}&\{7,14,21,28,35,42\}.
\end{align*}
The union of the first three lines is
\[
\Zfortynine\setminus Q
=\Zfortynine\setminus\{3,7,14,22,28,32,40,48\}.
\]
None of the first four lines meets \(Q\).  The last line
proves~\eqref{eq:PP-mod}, and direct inspection of the nine ordered pairs in
\(P\times P\) gives~\eqref{eq:diagonal-pairs}.
\end{proof}

We now define the infinite set
\begin{equation}\label{eq:S}
\begin{split}
S={}&D
\cup\{49q+b:q\in\Nzero,\ b\in B\}\\
&\cup\{49q+7:q\in Y_0\}\\
&\cup\{49q+14:q\in Y_1\}\\
&\cup\{49q+28:q\in Y_2\}.
\end{split}
\end{equation}
Every integer in~\eqref{eq:S} is positive.  Below \(98\), the five full
\(B\)-rails contribute
\[
\{3,22,32,40,48,52,71,81,89,97\}.
\]
Since
\[
Y_0\cap\{0,1\}=\{0\},
\quad
Y_1\cap\{0,1\}=\{1\},
\quad
Y_2\cap\{0,1\}=\{0,1\},
\]
the three controller rails contribute \(\{7,28,63,77\}\).  Together with
the seven elements of \(D\), this gives
\begin{equation}\label{eq:seed-prefix}
\begin{split}
S\cap\ints1{97}
=\{&1,2,3,5,7,13,22,27,28,32,36,40,47,48,52,\\
    &63,71,77,81,89,97\}.
\end{split}
\end{equation}

\section{The exact greedy recurrence}

\begin{theorem}\label{thm:recurrence}
For every \(t>97\),
\begin{equation}\label{eq:main-recurrence}
t\in S
\quad\Longleftrightarrow\quad
t\notin S+S.
\end{equation}
\end{theorem}

\begin{proof}
Write \(t=49q+r\), with \(0\leq r<49\).  Since \(t>97\),
\begin{equation}\label{eq:q-ge-two}
q\geq2.
\end{equation}

We first treat residues in \(Q=B\cup P\).  For \(t>97\), membership
\(t\in S\) can only come from one of the infinite rails
in~\eqref{eq:S}, and hence requires \(r\in Q\).  The six possible types of
pair sums are
\[
D+D,\quad D+B,\quad D+P,\quad B+B,\quad B+P,\quad P+P.
\]
A \(D+D\) sum is at most \(94<t\).  By~\eqref{eq:avoid-mod}, none of the next
four types can have residue in \(Q\).  By Lemma~\ref{lem:shield}, a \(P+P\)
sum can return to \(Q\) only through one of the three diagonal pairs
in~\eqref{eq:diagonal-pairs}.

It follows that, for the present \(q\geq2\), representation on the controller
rails is governed exactly as follows:
\begin{align*}
49q+14\in S+S
&\iff q\in Y_0+Y_0,\\
49q+28\in S+S
&\iff q\in Y_1+Y_1,\\
49q+7\in S+S
&\iff q\geq1\ \text{and}\ q-1\in Y_2+Y_2.
\end{align*}
Indeed, the corresponding representations are
\begin{align*}
(49u+7)+(49v+7)&=49(u+v)+14,\\
(49u+14)+(49v+14)&=49(u+v)+28,\\
(49u+28)+(49v+28)&=49(u+v+1)+7.
\end{align*}
The converse in each line follows from the exclusion of every other source
type.  Equations~\eqref{eq:gate01},~\eqref{eq:gate12},
and~\eqref{eq:gate20} now show that a controller-rail integer is
represented if
and only if it is omitted from \(S\).  No integer \(t>97\) on a full
\(B\)-rail is represented, by the preceding exhaustive six-type analysis.
Thus~\eqref{eq:main-recurrence} holds for every residue in \(Q\).
In particular,
if an omitted controller integer is considered explicitly, the relevant
complement identity supplies \(u,v\) in the source controller set, and the
three displayed equations give a representation by elements of \(S\).

It remains to treat \(r\notin Q\).  Such an integer is not in \(S\):
its residue excludes every infinite rail, and \(t>97\) excludes the finite
set \(D\).  Equation~\eqref{eq:cover-mod} supplies a residue representation
of one of the three
types \(B+B\), \(D+B\), or \(P+B\).

Suppose first that \(r=b+b'\) in \(\Zfortynine\), with \(b,b'\in B\).  For
the standard representatives, write
\(b+b'=49c+r\), where \(c\in\{0,1\}\).  Then
\[
t=b+\bigl(49(q-c)+b'\bigr).
\]
Both summands lie on full \(B\)-rails.  The second is positive, since
\(q\geq2\) and \(c\leq1\).  If instead \(r=d+b\) in
\(\Zfortynine\), with \(d\in D\) and \(b\in B\), write
\(d+b=49c+r\), where \(c\in\{0,1\}\).  Then
\[
t=d+\bigl(49(q-c)+b\bigr),
\]
using the fixed selected element \(d\) and a full-rail element.

Finally suppose that \(r=p+b\) in \(\Zfortynine\), with \(p\in P\) and
\(b\in B\).  The integers
\[
e_7=7,
\qquad
e_{14}=63=49+14,
\qquad
e_{28}=28
\]
belong to \(S\), because \(0\in Y_0\), \(1\in Y_1\), and \(0\in Y_2\).
Write \(e_p=49u_p+p\) and \(p+b=49c+r\).  Here
\(u_p,c\in\{0,1\}\), and
\[
t=e_p+\bigl(49(q-u_p-c)+b\bigr).
\]
By~\eqref{eq:q-ge-two}, the second quotient is nonnegative, so the second
summand is a positive element of a full \(B\)-rail.  This represents every
integer whose residue lies outside \(Q\).

All constructed summands are positive; since their sum is \(t\), each is
strictly smaller than \(t\).  We have therefore proved both directions
of~\eqref{eq:main-recurrence} for every \(t>97\).
\end{proof}

By~\eqref{eq:seed-prefix}, Theorem~\ref{thm:recurrence}, and
Lemma~\ref{lem:greedy}, the increasing enumeration of \(S\) is exactly the
greedy extension of the seed in Theorem~\ref{thm:main}.  All \(21\) seed
elements are installed before the first greedy test, which occurs after
\(97\).  No condition is imposed retroactively; for example, the seed is not
sum-free because \(1+1=2\).

\section{Nonperiodicity of the gaps}

\begin{lemma}\label{lem:gaps-imply-membership}
Let \(T=\{t_1<t_2<\cdots\}\subseteq\N\) be infinite.  If the consecutive
gaps \(t_{m+1}-t_m\) are eventually periodic, then the membership indicator
of \(T\) is ultimately periodic.
\end{lemma}

\begin{proof}
Suppose that the gaps have period \(p\geq1\) from index \(M\) onward.  Let
\[
h=\sum_{j=0}^{p-1}(t_{M+j+1}-t_{M+j})>0.
\]
Every block of \(p\) consecutive gaps in the periodic tail has the same sum,
so
\begin{equation}\label{eq:term-translation}
t_{m+p}=t_m+h
\qquad(m\geq M).
\end{equation}
The infinite set \(T\subseteq\N\) is unbounded.  Hence, given an integer
\(n\geq t_M\), there is a unique \(m\geq M\) such that
\(t_m\leq n<t_{m+1}\).  By~\eqref{eq:term-translation},
\[
t_{m+p}\leq n+h<t_{m+p+1},
\qquad
(n+h)-t_{m+p}=n-t_m.
\]
Thus \(n=t_m\) if and only if \(n+h=t_{m+p}\).  Equivalently,
\[
n\in T\iff n+h\in T
\qquad(n\geq t_M),
\]
which is ultimate periodicity of the membership indicator.
\end{proof}

\begin{proof}[Proof of Theorem~\ref{thm:main}]
The residue-\(7\) rail of \(S\) satisfies the exact identity
\begin{equation}\label{eq:rail-seven}
49q+7\in S
\quad\Longleftrightarrow\quad
q\in Y_0
\qquad(q\in\Nzero).
\end{equation}
Suppose, for contradiction, that membership in \(S\) were ultimately
periodic with some period \(h\geq1\).  Every positive multiple of an ultimate
period is again an ultimate period, so \(49h\) is an ultimate period of \(S\).
For every sufficiently large \(q\), equations~\eqref{eq:rail-seven} and the
\(49h\)-periodicity of \(S\) would therefore give
\[
q\in Y_0
\iff 49q+7\in S
\iff 49(q+h)+7\in S
\iff q+h\in Y_0.
\]
This contradicts Lemma~\ref{lem:Y0-aperiodic}.  Hence the membership
indicator of \(S\) is not ultimately periodic.  By the contrapositive of
Lemma~\ref{lem:gaps-imply-membership}, the consecutive gaps in the increasing
enumeration of \(S\) are not eventually periodic.  This proves the asserted
quantified statement for every \(M,p\geq1\).
\end{proof}

\end{document}
